Metal-poor stars preserve a fossil record of the early Milky Way, but extinction and crowding make them difficult to identify in the Galactic midplane. We propose a 60 deg² pathfinder survey with Blanco/NEWFIRM to identify metal-poor luminous giants for spectroscopic follow-up by After-SDSS-V’s Hidden Galaxy Explorer (HGE), which will begin mapping the dust-obscured far side of the Milky Way in mid-2027. Using the H1 and H2 intermediate-band filters, we will survey 40° < l < 60° and |b| < 1.5°. Synthetic photometry predicts that H1−H2 is sensitive to metallicity along the red-giant branch, while existing J−Ks photometry constrains temperature and reddening. The footprint is expected to contain approximately 13,000 giants with [Fe/H] < −1, including about 1,300 with [Fe/H] < −2. In two nights, we will image 306 pointings in both filters, reaching S/N ≈ 140 per band at H = 15.5 and σ(H1−H2) ≈ 0.011 mag. Existing APOGEE spectroscopy will calibrate the selection at the bright end, while HGE spectra will validate it at fainter magnitudes. The survey will provide HGE with a ranked catalog of high-priority metal-poor candidates and test the strategy needed for future expansion across its broader footprint.